\section{Kronecker coefficients}\label{sec:kronecker}

Let $W=\C^n$ with basis $e_1,\dots,e_n$, and let $R_n=\C[W\otimes W\otimes W]
^{SL_n^3}$, graded so that $(R_n)_s$ is the space of invariants of degree $ns$
(invariants occur only in degrees divisible by $n$). By Schur--Weyl duality
$\dim(R_n)_s=g((s^n),(s^n),(s^n))=k_n(s)$, so $K_n=\Hilb_{R_n}$. Let $\tau$ be
the cyclic permutation of the three tensor factors; it acts on $R_n$ as a
graded algebra automorphism of order three.

Let $\mathcal U_n=W\oplus\mathbb S_{(2,1)}W$, give its coordinate ring the
grading in which linear functions on $W$ have weight one and those on
$\mathbb S_{(2,1)}W$ weight three, and let $b_n(s)$ be the dimension of the
$SL_n$-invariants of weighted degree $ns$ and $B_n(t)=\sum_sb_n(s)t^s$. (All
invariants have weighted degree divisible by $n$, as the scalars $\mu_n$ act
on weighted degree $D$ by $\zeta^D$.) Then $B_n=\Hilb_{T_n}$ with
$T_n=\C[\mathcal U_n]^{SL_n}$ graded by $s$.

\begin{theorem}\label{thm:K023}
For every $n\ge2$ and $s\ge0$,
\[
 \tr\bigl(\tau\mid(R_n)_s\bigr)=b_n(s),\qquad k_n(s)\equiv b_n(s)\pmod3 ,
\]
and $k_n(s)-b_n(s)$ is a nonnegative multiple of three. Hence
$K_n(t)\equiv B_n(t)\pmod{3\Z[[t]]}$.
\end{theorem}

\begin{proof}
We may compute with polynomial representations of $W$ instead of $W^*$; the
simultaneous contragredient preserves the dimensions of invariant spaces and
commutes with permutations of the factors. In degree $d=ns$, Schur--Weyl
duality gives, compatibly with $\tau$,
\[
 (R_n)_s\simeq\bigl([\lambda]\otimes[\lambda]\otimes[\lambda]\bigr)^{S_d},
 \qquad\lambda=(s^n),
\]
since $\Sym^d(W\otimes W\otimes W)$ is the diagonal $S_d$-invariant part of
$(W^{\otimes d})^{\otimes3}$ and $(W^{\otimes d})^{SL_n}$ is the Specht module
$[\lambda]$. For an endomorphism $A$ of a vector space, the trace of the cyclic
permutation composed with $A^{\otimes3}$ is $\tr(A^3)$ (contract the cyclic
indices in a basis). Averaging over $S_d$,
\[
 \tr\bigl(\tau\mid(R_n)_s\bigr)=\frac1{d!}\sum_{\pi\in S_d}\chi_\lambda(\pi^3).
\]
By the Frobenius formula $p_\mu=\sum_\lambda\chi_\lambda(\mu)s_\lambda$, this
average is the coefficient of $s_\lambda z^d$ in
$\sum_d\frac{z^d}{d!}\sum_{\pi\in S_d}p_{{\rm type}(\pi^3)}$. Cubing a
permutation keeps a $j$-cycle with $3\nmid j$ and splits a $3r$-cycle into three
$r$-cycles, so by the exponential formula the series is
\[
 \exp\Bigl(\sum_{3\nmid j}\frac{z^jp_j}j+\sum_{r\ge1}\frac{z^{3r}p_r^3}{3r}\Bigr)
 =\exp\Bigl(\sum_{r\ge1}\frac{z^rp_r}r+\sum_{r\ge1}\frac{z^{3r}}r\,
 p_r[s_{(2,1)}]\Bigr),
\]
using $p_r[s_{(2,1)}]=(p_r^3-p_{3r})/3$, which follows from
$s_{(2,1)}=(p_1^3-p_3)/3$. The right side is the character of the symmetric
algebra on $W\oplus\mathbb S_{(2,1)}W$ with weights one and three
\cite[Chapter I, \S\S7--8]{Macdonald}; specializing to $n$ variables and taking
the coefficient of $s_{(s^n)}$ gives $b_n(s)$. Finally, $\tau$ is defined over
$\Q$ on $(R_n)_s$, and an operator of order three on a rational vector space has
characteristic polynomial $(X-1)^a(X^2+X+1)^b$, dimension $a+2b$ and trace
$a-b$; the difference is $3b\ge0$.
\end{proof}

We realize $\mathbb S_{(2,1)}W$ as the degree-three part of the free Lie
algebra on $W$ inside $W^{\otimes3}$, with $[a,b]=ab-ba$: the map
$W\otimes\Lambda^2W\to W^{\otimes3}$, $a\otimes(b\wedge c)\mapsto[a,[b,c]]$,
has kernel $\Lambda^3W$ by the Jacobi identity, and its image is
$\mathbb S_{(2,1)}W$. Its weights are $2e_i+e_j$ ($i\ne j$), with multiplicity
one, and $e_i+e_j+e_k$ ($i,j,k$ distinct), with multiplicity two. Put
\[
 z_n=(e_n,F_n)\in\mathcal U_n,\qquad F_n=\sum_{i=1}^{n-1}[e_i,[e_i,e_{i+1}]],
 \qquad M=2^{n-1}.
\]

\begin{lemma}\label{lem:liezero}
For every $n\ge2$, the stabilizer of $z_n$ in $\mathfrak{gl}_n$ is zero, so the
stabilizer of $z_n$ in $GL_n$ is finite.
\end{lemma}

\begin{proof}
Write $Xe_j=\sum_ax_{aj}e_a$ and suppose $Xz_n=0$; then $Xe_n=0$, and $X$ acts
on $F_n$ as a derivation of the brackets. For $r\ne s$, the weight $2e_r+e_s$
is spanned by $L_{rs}=[e_r,[e_r,e_s]]$, and its coefficient in $XF_n$ is
\[
 \mathbf1_{r<n}x_{s,r+1}+2\cdot\mathbf1_{s=r+1}x_{rr}-\mathbf1_{s=r-1}x_{r,r-1},
\]
the last term coming from $[e_r,[e_{r-1},e_r]]=-L_{r,r-1}$. Taking $s=r+1$ gives
$2x_{rr}+x_{r+1,r+1}=0$, and $x_{nn}=0$, so the diagonal vanishes. For column
$j\ge2$ the only possibly nonzero entries are then $x_{j-1,j}$ and
$x_{j-2,j}=x_{j-1,j-2}$ ($j\ge3$); for $j\ge4$ the latter is an entry below the
diagonal in column $j-2\ge2$, hence zero. So besides the first column and the
superdiagonal only $x_{13}=x_{21}$ ($n\ge3$) can be nonzero. Put
$y_j=x_{j-1,j}$; $Xe_n=0$ gives $y_n=0$. If $n=2$, $Xe_2=0$ also kills $x_{12}$,
and the coefficient of $L_{21}$ gives $x_{21}=0$. Let $n\ge3$.

For $a\ge4$, the weight $e_1+e_2+e_a$ receives only the term of index one,
with coefficient $x_{a1}$ times the nonzero vector
$[e_a,[e_1,e_2]]+[e_1,[e_a,e_2]]$; so $x_{a1}=0$. For $3\le j\le n-1$, the
weight $e_{j-1}+e_j+e_{j+1}$ receives only the two adjacent terms, with
coefficients $y_j$ and $x_{j+1,j-1}=0$, the former multiplying the nonzero
vector $[e_{j-1},[e_j,e_{j+1}]]+[e_j,[e_{j-1},e_{j+1}]]$; so $y_j=0$. In the
weight $e_1+e_2+e_3$, with the basis $A=[e_1,[e_2,e_3]]$,
$B=[e_2,[e_1,e_3]]$ and $[e_3,[e_1,e_2]]=-A+B$, the contribution is
$y_2(A+B)+x_{31}(-2A+B)$; so $y_2-2x_{31}=y_2+x_{31}=0$ and
$x_{31}=y_2=0$. Finally $x_{13}=x_{21}$ vanishes by $Xe_3=0$ if $n=3$, and for
$n\ge4$ because the weight $e_1+e_3+e_4$ receives only the derivative of the
repeated $e_3$ in the term of index three, with coefficient $x_{13}$ times a
nonzero vector. So $X=0$. An algebraic group in characteristic zero has the
dimension of its Lie algebra.
\end{proof}

\begin{theorem}\label{thm:K024}
Let $n\ge2$, $M=2^{n-1}$, $\zeta$ a primitive $M$-th root of unity, and
$X_n=\Spec T_n$.
\begin{enumerate}
\item The $SL_n$-orbit of $z_n$ is closed; its stabilizer in $GL_n$ is cyclic of
order $M$, generated by
\begin{equation}\label{eq:chainelement}
 g=\diag(\zeta^{w_1},\dots,\zeta^{w_n}),\qquad w_i=(-2)^{i-1}\bmod M ,
\end{equation}
on which $\det$ is faithful; its stabilizer in $SL_n$ is trivial.
\item The nonzero closed points of $X_n$ fixed by the grading group $\mu_M$
form one orbit of the grading group $\Gm$, the orbit of the image $p_n$ of
$z_n$, whose stabilizer in $\Gm$ is exactly $\mu_M$. $X_n$ is smooth at $p_n$,
and the fixed scheme $X_n^{\zeta}$ is smooth of dimension one there.
\item $B_n$ has a simple pole at every primitive $M$-th root of unity $\xi$,
with $\lim_{t\to\xi}(1-t^M)B_n(t)$ a unit in $\Z[\xi,1/2]$.
\item If the two-part of an integer $d$ is at least $M$ and $d\ne M$, then
$B_n$ is regular at every primitive $d$-th root of unity; in particular at the
primitive $M3^r$-th roots for $r\ge1$.
\end{enumerate}
\end{theorem}

\begin{proof}
\emph{Closed orbit.} For $1\le i<n$ put $q_i=(1-(-1/2)^i)/3$ and $q_n=1-q_{n-1}$;
all are positive, and with $a_i=2e_i+e_{i+1}$ ($i<n$), $a_n=e_n$, one has
$\sum_iq_ia_i=\one$ ($2q_1=1$, $q_{i-1}+2q_i=1$, $q_{n-1}+q_n=1$). The words
$[e_i,[e_i,e_{i+1}]]$ have squared norm six in the Hermitian metric of
$W^{\otimes3}$; give the word $i$ the coefficient $\sqrt{q_i/6}$ and $e_n$ the
coefficient $\sqrt{q_n}$. The diagonal part of the moment map of $SU_n$ is then
$\one$ minus its trace part, which vanishes; its off-diagonal part vanishes
because no two of the weights $a_1,\dots,a_n$ in the same summand differ by a
root $e_a-e_b$, and the summands are orthogonal and preserved. So this vector
is a zero of the moment map, and its $SL_n$-orbit is closed by the
Kempf--Ness criterion \cite[Theorem 4, p.~1]{Losev}. The weights $a_i$ are
linearly independent (triangular, determinant $2^{n-1}$), so the diagonal torus
acts transitively on vectors with these weights and nonzero coefficients; it
maps our vector to $z_n$, and the orbit of $z_n$ is closed as well.

\emph{Stabilizer.} The weights $w_i$ are distinct modulo $M$, $w_n=0$ and
$2w_i+w_{i+1}\equiv0$, so $g$ fixes $z_n$, and $\det g=\zeta^{\sum w_i}$ with
$\sum_iw_i$ odd. Let $L=\Stab_{GL_n}(z_n)$, finite by Lemma~\ref{lem:liezero}.
As $g$ has simple spectrum, its centralizer in $L$ is diagonal; a diagonal
element of $L$ satisfies $d_n=1$ and $d_i^2d_{i+1}=1$, so the centralizer is
$\langle g\rangle$, of order $M$. Column orthogonality gives
$\sum_{\rho\in{\rm Irr}(L)}|\rho(g)|^2=M$; the powers $\det^k$, $0\le k<M$, are
distinct linear characters (as $\det g$ has order $M$) and exhaust this sum. So
every linear character of $L$ is a power of $\det$, and every nonlinear
irreducible character vanishes at $g$. Let $K=\ker(\det|_L)$. A representation
of $g$ has trace zero exactly when the multiplicities of $\eta$ and $-\eta$
agree for every eigenvalue $\eta$ (use the basis $1,\zeta,\dots,\zeta^{M/2-1}$
of $\Q(\zeta)$ and $\zeta^{M/2}=-1$). The spectrum \eqref{eq:chainelement} has
exactly one pair $\{\eta,-\eta\}$, on $P=\langle e_{n-1},e_n\rangle$, since
the $w_i$ with $i\le n-2$ have distinct two-adic valuations less than $n-2$.
So every nonlinear constituent of $W$ lies in $P$; the linear constituents are
powers of $\det$, on which $K$ acts trivially; $P$ is $L$-stable, $K$ acts on
it with determinant one and fixes $e_n$, and a finite-order element of
$SL_2(\C)$ fixing a nonzero vector is the identity. So $K=1$ and
$L=\langle g\rangle$; the $SL_n$-stabilizer is trivial.

\emph{The grading.} The scalar $\lambda I\in GL_n$ acts on $\mathcal U_n$ as the
weight grading, and on $(T_n)_s$ by $\lambda^{\pm ns}=(\det\lambda I)^{\pm s}$;
so each $B\in GL_n$ acts on $X_n$ as the grading element $(\det B)^{\pm1}$.
Hence the grading stabilizer of $p_n$ is $\det(L)=\mu_M$, and every point of
$\Gm\cdot p_n$ is fixed by $\mu_M$.

\emph{Fixed points.} Let $p\ne0$ be a closed point of $X_n$ fixed by a grading
element of order $2^a\ge M$ and $y$ in its closed orbit. By
Lemma~\ref{lem:lift} (with $\gamma=\pm1$) some diagonal $A$ of two-power order
fixes $y$ with $\ord(\det A)=2^a$. In an eigenbasis of $A$ consider the weights
of the nonzero components of $y$: linear weights $e_i$, and weights $2e_i+e_j$,
$e_i+e_j+e_k$ of $\mathbb S_{(2,1)}W$. Replace each linear weight $e_i$ by
$3e_i$; this preserves the rational span and, for eigenvalues of two-power
order, the condition of being fixed ($\alpha_i=1\iff\alpha_i^3=1$). The span
condition holds by the argument of Lemma~\ref{lem:span}, applied to a
positive-degree invariant nonzero at $y$. The resulting columns are the support
of a cubic form fixed by $A$, and Theorem~\ref{thm:fullcubicsupport} gives
$2^a\le M$; so $a=n-1$, and the support is the chain \eqref{eq:chainsupport}
after a permutation. The cube $3e_n$ comes from the linear coordinate $e_n$,
since $\mathbb S_{(2,1)}W$ has no weight $3e_i$, and the other columns from the
one-dimensional weight spaces of $2e_i+e_{i+1}$. So $y$ is diagonally
equivalent to $z_n$, and $p\in\Gm\cdot p_n$.

\emph{Smoothness and the fixed tangent space.} Since the closed orbit of $z_n$
has trivial stabilizer in $SL_n$, Luna's slice theorem
\cite[Section III.1, Th\'eor\`eme du slice \'etale and Remarque 1$^\circ$, pp.~97--98]{Luna} shows that $X_n$ is smooth at $p_n$ with
tangent space $\mathcal U_n/\mathfrak{sl}_n\cdot z_n$, equivariantly for $g$.
The fixed space of $g$ on $\mathcal U_n$ has dimension $n$. Recall
$v_2(w_i)=i-1$ for $i<n$ and $w_n=0$. A linear weight $e_i$ is fixed only for
$i=n$. A weight $2e_i+e_j$ is fixed only for $j=i+1$: if $j=n$, then
$2w_i\equiv0$ forces $i=n-1$; if $i=n$, then $w_j\not\equiv0$; if $i,j<n$ and
$j\ne i+1$, the valuations $i$ and $j-1$ of $2w_i$ and $w_j$ differ, so the sum
has valuation at most $n-2$; and $2w_i+w_{i+1}=0$. A weight $e_i+e_j+e_k$ is
never fixed, since among its nonzero terms the least valuation occurs once.
As $g$
has simple spectrum, $\dim\mathfrak{sl}_n^g=n-1$. Taking fixed vectors is exact
for a finite group in characteristic zero, so $\dim(T_{p_n}X_n)^g=1$, and as in
the proof of Proposition~\ref{prop:fixedcomponentpole} the fixed scheme is
smooth of dimension one at $p_n$. Translating by the grading group, the same
holds along $\Gm\cdot p_n$.

\emph{The pole.} Let $Y$ be the closure of $\Gm\cdot p_n$. Its only other point
is the origin, since the grading is positive; so the support of
$X_n^{\zeta}$ is $Y$, a curve, and $X_n$ is smooth at its generic point. By
Lemma~\ref{lem:orbitclosure}, $e_Y=1$. Proposition~\ref{prop:fixedcomponentpole}
gives a simple pole at $\zeta$ with $\lim(1-t^M)B_n=1/\prod_j(1-\nu_j)$, where
the $\nu_j$ are the eigenvalues of $\zeta$ on the conormal space, all
different from one and all $M$-th roots of unity. For an $M$-th root of unity
$\nu\ne1$ of order $2^b$, $(1-\nu)\sum_{j<2^{b-1}}\nu^j=1-\nu^{2^{b-1}}=2$, so
$1-\nu$ is a unit in $\Z[\zeta,1/2]$. The same holds at every primitive
$M$-th root by Galois conjugation, $B_n$ having rational coefficients.

\emph{No other poles of large two-part.} If the two-part of $d$ is at least $M$
and $d\ne M$, a nonzero point fixed by a primitive $d$-th root would be fixed by
its two-primary part, of order at least $M$, hence lie on $\Gm\cdot p_n$, whose
grading stabilizer is $\mu_M\not\ni$ the primitive $d$-th roots. So the fixed
locus is the origin, and Lemma~\ref{lem:isolated} applies.
\end{proof}

\begin{theorem}\label{thm:kronecker}
For every $n\ge2$ there is $r\ge0$ such that $K_n(t)=\sum_sg((s^n),(s^n),
(s^n))t^s$ has a pole at every primitive root of unity of order $2^{n-1}3^r$.
Every expression $K_n(t)=P(t)/\prod_i(1-t^{a_i})^{b_i}$ with $P\in\Q[t]$ and
positive integers $a_i,b_i$ has some $a_i$ divisible by $2^{n-1}3^r$; in
particular every positive form of $K_n$ has modular index at least $2^{n-1}$.
\end{theorem}

\begin{proof}
Both $K_n$ and $B_n$ are Hilbert series of finitely generated positively graded
domains, so they are rational with poles at roots of unity, with integral
coefficients. By Theorem~\ref{thm:K023}, $K_n\equiv B_n\pmod 3$. By
Theorem~\ref{thm:K024}, $B_n$ has a simple pole at each primitive $M$-th root
$\zeta$, $M=2^{n-1}$, with limit a unit in $\Z[\zeta,1/2]$, hence a unit at
every prime above $3$, and no pole at the primitive $M3^r$-th roots, $r\ge1$.
Theorem~\ref{thm:transfer} with $p=3$, $F=K_n$, $G=B_n$ applies.
\end{proof}

\begin{corollary}\label{cor:kroneckerPH3}
The existence clause of PH3 for Kronecker coefficients \cite[Hypothesis
3.3.4]{GCT6} fails: for the inputs $((1^n),(1^n),(1^n))$, with
$\rank(x)=n^2+5n$, every positive rational form of the generating function has
modular index at least $2^{n-1}$. The relaxed PH3 of \cite[Appendix,
p.~20]{BOR}, which bounds the modular index of the minimal positive form by a
polynomial in the largest height $h=n$, fails as well.
\end{corollary}

\begin{proof}
The minimal positive form is a positive form, and $2^{n-1}$ exceeds every
polynomial in $n$, hence in $n^2+5n$, for large $n$.
\end{proof}

